\section{Introduction}
\label{sec:intro}

Two observations about the current state of language modelling sit in quiet tension. The first is that the abilities we associate with reasoning (multi-step arithmetic, long-context inference, writing working code) appear, by a range of measurements, only once standard English-centric transformers exceed $10^{10}$ parameters \citep{wei2022}. The second is that the parameter counts a consumer device can actually hold (smartphones, laptops with integrated GPUs, embedded devices) sit well below this threshold, even after aggressive compression to 4 bits per weight or lower (roughly $5 \times 10^{8}$ to $2 \times 10^{9}$ parameters on current mobile chips). The gap between these two figures is why reasoning-capable language models live, at present, as cloud services rather than on the device in the user's hand.

Most treatments of this gap quietly frame it as a problem of scale: a matter of model compression, weight quantization, knowledge distillation from a larger teacher, or architectural cleverness that will eventually close the distance between what a model can do and what a device can hold. What that framing leaves unexamined is that the figures on both sides of the gap are fit to English. Transformer scaling laws \citep{kaplan2020, hoffmann2022} are calibrated on predominantly English corpora. The benchmarks tracking the onset of reasoning likewise inherit English as their default experimental medium. Whether the same parameter thresholds hold for languages whose grammar is organised quite differently from English (those whose word forms pack more grammatical information into a single token, or whose internal structure is more easily read off the surface form) is not a settled question; for the most part it is not even an asked question.

Responses to the gap between edge and reasoning fall along three axes the field rarely names in one breath. The first, and most heavily funded, is \emph{brute-force hardware scaling}: larger datacentre GPUs (Nvidia's H- and B-series), more high-bandwidth memory, more interconnect, more electrical power, on the bet that capability per dollar will keep bending favourably with scale. The second, more speculative, is \emph{substrate redefinition}: companies such as Extropic\footnote{\url{https://extropic.ai}} are building \emph{thermodynamic sampling units}, chips that are inherently probabilistic, arguing that modern AI inference is at heart a sampling workload whose deterministic emulation on transistor silicon is physically wasteful. The third, which this paper develops, sits one layer above both. Rather than adding hardware or rewriting its physics, it aligns the tokenisation layer to the grammatical structure of the language being modelled, freeing up compute on the existing GPU substrate that is currently spent on learning grammar from distributional statistics. The three routes are not mutually exclusive, but they cut the cost problem very differently. Nvidia spends more energy to train the same languages faster; Extropic proposes spending much less energy under a different computational paradigm; morphology-aware tokenisation spends nothing extra and recovers compute by changing what the model has to learn in the first place.

This paper develops the possibility that the parameter threshold for a given capability depends, in a measurable way, on the language. More precisely: a transformer trained under a tokenisation that respects a language's morphological structure, where \emph{grammatical feature bundles} (the parsed grammatical analysis of each word) are handed to the model as explicit categorical inputs rather than reconstructed from the co-occurrence of byte-pair fragments, is spared a portion of the parameter budget its byte-pair counterpart must spend on rebuilding that structure from scratch.

A \emph{grammatical feature bundle}, in the sense used throughout this paper, is the complete set of grammatical category and value pairs that a single word form realises at once. The English \emph{researchers} has the bundle $\{\text{NOUN}, \text{number}=\text{plural}\}$; the Turkish \emph{evlerinizden} carries $\{\text{NOUN}, \text{number}=\text{plural}, \text{possessor}=\text{2-person plural}, \text{case}=\text{ablative}\}$; the Arabic \ar{اسْتَخْدَمَ} carries $\{\text{VERB}, \text{form}=\text{X}\}$. Bundles are the unit of grammatical information the framework quantifies. They are what a morphology-aware tokeniser encodes outright, what a byte-pair tokeniser leaves the model to recover from distributional statistics, and what the two quantities introduced next, $H(L)$ and $\rho(L)$, respectively measure the Shannon entropy and the structural recoverability of.

The size of the rebate the framework argues for is set by two structural quantities of the language: the Shannon entropy $H(L)$ of its grammatical feature bundles per word form, that is, how much grammatical information the average word actually carries, and the coefficient $\rho(L)$, which measures what fraction of that information is recoverable just by looking at the surface form. Both quantities can be computed from a morphological grammar engine on its own, before any training run is launched. The framework predicts that the product $\rho(L) \cdot H(L)$ orders the observed parameter rebate across languages, and, under a stated linearity assumption, that a single cross-lingual constant $\alpha$ turns the product into a quantitative prediction.

The paper reports a six-model Phase~1 experiment (English, Arabic, and Turkish, each under byte-pair and morphology-aligned tokenisation, identical in every other respect) testing both forms of this prediction. The ordinal prediction holds: the observed efficiency gain is ordered by $\rho \cdot H$ across the three languages, with English smallest and Turkish largest, in exact agreement with the framework's forward calculation from the grammar engines. The magnitude prediction does not hold: the implied per-bit coefficient $\alpha$ varies by roughly forty-fold across the three languages under any reasonable parameterisation of the Hoffmann scaling law. The shape of the deviation is not random: the implied $\alpha$ itself grows with $\rho \cdot H$. Taken at face value, Phase~1 suggests that structurally recoverable grammar buys parameter savings at an accelerating rate, a pattern we name the \textbf{Agglutinative Compounding Effect}, and whose reading the paper leaves open between two scientifically distinct possibilities.

The paper makes four contributions at once. The first is a framework in which a language's morphology maps to a quantitative prediction about transformer parameter requirements, with inputs that can be computed before any training begins. The second is a Phase~1 ordinal confirmation of the framework's prediction across three typologically distant languages. The third is a named empirical phenomenon, the \textbf{Agglutinative Compounding Effect}, which is consistent with two mutually exclusive readings: one under which morphology acts as a compounding lever on model capacity, and one under which the standard scaling laws require a morphological correction term.

The fourth is a companion phenomenon, the \textbf{Templatic Lexical Surplus}, motivated by two converging observations. The first was mechanical: the Arabic grammar engine enumerated roughly 270 distinct feature bundles, against Turkish's 63 and English's 23, a spread wide enough to raise the question of what that count actually \emph{means} for the density of information packed into a word. Does Arabic genuinely carry more grammatical information per word, or is its bundle taxonomy simply more combinatorial? The second observation came from outside computational linguistics. \citet{wei2023native}, imaging 94 native Arabic and German speakers at the Max Planck Institute for Human Cognitive and Brain Sciences, found that Arabic speakers show measurably stronger connectivity in left temporo-parietal regions (the semantic language areas) and stronger inter-hemispheric connections through the posterior corpus callosum. The authors attribute this to the processing demands of Arabic's root-and-pattern morphology. Both lines suggested that Arabic's true per-word information load is larger than its surface behaviour reveals, and that the surface-only $\rho(L)$ of our main framework therefore systematically under-prices Arabic's rebate. The two-tier extension tests this directly. A second input stream, exposing the root consonants themselves as a parallel channel feeding the model alongside the surface form and the feature bundle, recovers an additional rebate of 0.082 nats at rung~1, holding remarkably stable across the three-rung ladder ($+0.077$ and $+0.058$ at rungs~2 and 3) even as the first-tier advantage itself inverts. The measured surplus only captures root identity. Arabic's \emph{awz\=an} (the templatic patterns themselves) additionally encode semantic function (agent, patient, instrument, \emph{ma\d{s}dar}, \emph{\d{s}ifa mushabbaha}, the Form-II causative, the Form-X seeking pattern, and similar first-class semantic classes), which the current engine collapses into coarse templatic categories. Isolating the wazn's semantic-function contribution as its own measurement, through a natural third-tier extension, is scoped as follow-up work in \S\ref{sec:conclusion}. The $0.082 / 0.077 / 0.058$ figures therefore understate the true Templatic Lexical Surplus: they count only the rebate attributable to root identity, and not the rebate the wazn's semantic function would additionally contribute if exposed. The full surplus is at least these numbers, and in all likelihood larger.

The paper takes no side between the two readings of the Compounding Effect, and specifies the follow-up experiments that would decide between them.

Edge-hostable reasoning, that is, reasoning-capable language models that fit inside smartphone-class parameter budgets, is the motivation for asking whether grammar is a capacity lever, and it is the setting in which the question's practical stakes are highest. It is not claimed as a result of this paper. It is framed as a conditional. \emph{If} the Agglutinative Compounding Effect is a genuine property of morphology-aware language modelling rather than an artefact of the Phase~1 regime, then edge-hostable reasoning becomes a plausible near-term prospect for morphologically rich languages, at parameter counts English does not currently reach. Whether that conditional holds is the empirical question Phase~1 could not settle on three data points, and that the research programme set out in \S\ref{sec:conclusion} is designed to answer.

The remainder of the paper is organised as follows. \S\ref{sec:related} places the work inside five neighbourhoods of prior art: morphological typology, tokenisation in NLP, transformer scaling laws, the framing of capability emergence, and the Green AI, linguistic equity, and edge ML literature on cost. \S\ref{sec:framework} develops the framework formally. \S\ref{sec:cases} walks through English, Arabic, and Turkish morphology, with the grammar-engine decompositions that supply the framework's inputs. \S\ref{sec:results} reports Phase~1, the ordinal validation, the magnitude divergence, and the naming of the Agglutinative Compounding Effect. \S\ref{sec:discussion} discusses the consequences of each interpretation, together with the limitations that bound confidence in the observation. \S\ref{sec:conclusion} concludes and specifies the discriminating research programme.
